\section{技术实现}
\label{sec:method}

本节详述各流水线阶段的技术实现。

\subsection{多源异构文献检索}
\label{sec:method-search}

\texttt{SearchAgent}通过以下步骤实现多数据库文献检索：

\textbf{查询构建。}给定研究主题，Agent首先调用LLM生成结构化搜索策略。提示词强制要求简洁的关键词组合（每条查询3--5个词）、数据库特定语法，以及通过领域限定词消除歧义。若LLM未能返回有效查询，回退策略将从研究主题中生成关键词两两和三三组合。

\textbf{多数据库搜索。}系统集成五个学术数据库，均无需API密钥即可访问：

\begin{table}[htbp]
    \centering
    \caption{支持的学术数据库}
    \label{tbl:databases}
    \begin{tabular}{llcc}
    \toprule
    数据库 & 领域 & 摘要 & 需要API Key \\
    \midrule
    arXiv       & 物理、计算机、数学     & 有 & 否 \\
    PubMed      & 医学、生物学           & 有 & 否 \\
    DBLP        & 计算机科学             & 无 & 否 \\
    Europe PMC  & 生命科学、生物医学     & 有 & 否 \\
    OpenAlex    & 跨学科                 & 无 & 否 \\
    \bottomrule
    \end{tabular}
\end{table}

每个数据库实施独立的速率限制，其中arXiv要求查询间3秒延迟以避免被封禁。若某次arXiv查询返回零结果，则跳过后续arXiv查询。

\textbf{去重与标准化。}检索到的论文基于标题标准化比较进行去重。每篇论文标准化为\texttt{LiteratureRecord}数据结构，包含标题、作者、摘要、年份、来源、DOI、期刊/会议和引用数字段。

\subsection{三级相关性筛选}
\label{sec:method-screen}

\texttt{ScreenAgent}实现三级筛选机制，在吞吐量和精度之间取得平衡。

\textbf{第一级：基于嵌入的余弦相似度。}
对每篇论文，计算其嵌入（标题+摘要）与每个RQ嵌入（问题+关键词）之间的余弦相似度。设$\mathbf{p}_i \in \mathbb{R}^d$为论文$i$的嵌入，$\mathbf{r}_j \in \mathbb{R}^d$为RQ $j$的嵌入，论文$i$的相关性分数为：
\begin{equation}
    s_i = \max_{j} \frac{\mathbf{p}_i \cdot \mathbf{r}_j}{\|\mathbf{p}_i\| \cdot \|\mathbf{r}_j\|}
    \label{eq:cosine-similarity}
\end{equation}
通过矩阵乘法$\mathbf{P} \mathbf{R}^\top$一次性计算，其中$\mathbf{P} \in \mathbb{R}^{n \times d}$和$\mathbf{R} \in \mathbb{R}^{m \times d}$分别为论文和RQ的嵌入矩阵。

\textbf{第二级：基于阈值的分类。}
论文根据相似度分数分为三个区域：
\begin{itemize}
    \item \textbf{自动通过}（$s_i \geq \theta_{\max}$）：直接标记为相关
    \item \textbf{自动拒绝}（$s_i < \theta_{\min}$）：直接标记为不相关
    \item \textbf{边界区域}（$\theta_{\min} \leq s_i < \theta_{\max}$）：送入LLM判定
\end{itemize}
其中$\theta_{\min} = 0.68$、$\theta_{\max} = 0.80$为可配置阈值。

\textbf{第三级：LLM辅助边界判定。}
对边界区域论文，LLM从三个维度进行加权评分：
\begin{equation}
    w_i = 0.5 \cdot \text{topic}_i + 0.3 \cdot \text{method}_i + 0.2 \cdot \text{timeliness}_i
    \label{eq:weighted-score}
\end{equation}
当且仅当$w_i \geq 3.0$时，论文被标记为相关。边界论文按相似度分数排序，限制最多150次LLM调用以控制成本。

\textbf{嵌入模型选择。}系统自动检测论文内容的语言，选择合适的嵌入模型：
\begin{itemize}
    \item 中文文本：BGE-small-zh-v1.5（768维，本地GPU推理）
    \item 英文文本：all-MiniLM-L6-v2（384维，本地GPU推理）
    \item 回退方案：远程GLM embedding-3 API（1024维）
\end{itemize}

\subsection{GPU加速语义聚类}
\label{sec:method-cluster}

\texttt{ClusterAgent}对相关论文进行多步骤语义聚类。

\textbf{降维。}
采用两步降维策略。给定嵌入矩阵$\mathbf{E} \in \mathbb{R}^{n \times d}$：
\begin{enumerate}
    \item \textbf{第一步：}PCA将维度从$d$降至50，GPU可用时使用PyTorch CUDA SVD分解：
    \begin{equation}
        \mathbf{E}_{50} = \mathbf{U}_{(:, 1:50)} \cdot \text{diag}(\mathbf{S}_{1:50})
        \label{eq:pca-reduction}
    \end{equation}
    其中$\mathbf{E} - \bar{\mathbf{E}} = \mathbf{U} \mathbf{S} \mathbf{V}^\top$为SVD分解。
    \item \textbf{第二步：}t-SNE从50维降至2维用于可视化，困惑度根据样本量自适应：
    \begin{equation}
        \text{perplexity} = \min(30, \max(5, n / 10))
        \label{eq:perplexity}
    \end{equation}
\end{enumerate}
对于小数据集（$n < 10$），直接使用PCA降至2维，跳过t-SNE。

\textbf{HDBSCAN聚类。}
HDBSCAN算法在降维空间中聚类论文，自适应参数：
\begin{equation}
    \text{min\_cluster\_size} = \max(2, n / 10)
    \label{eq:min-cluster}
\end{equation}
无法分配到任何簇的论文标记为噪声（$-1$）。当HDBSCAN不可用时，KMeans作为回退方案。

\textbf{聚类标签生成。}
对每个聚类，LLM生成描述性标签、核心主题、代表性论文、子主题，以及该簇论文共享的核心公式或数学框架和底层物理/数学原理。

\textbf{领域一致性过滤。}
通过比较聚类论文标题与研究主题关键词的重叠率进行过滤。主题关键词匹配率低于30\%的聚类被移除。

\subsection{聚焦原理的报告生成}
\label{sec:method-summary}

\texttt{SummaryAgent}通过精心设计的提示词工程流水线生成结构化综述报告。

\textbf{两阶段生成。}
最终报告通过两次LLM调用生成，以绕过输出token限制：
\begin{itemize}
    \item \textbf{第一部分：}引言 + 核心技术原理（按聚类分析）
    \item \textbf{第二部分：}应用分析 + 技术演进 + 理论极限 + 突破路径 + 结论
\end{itemize}

\textbf{反幻觉提示词设计。}
提示词在多个层级强制执行严格规则：

\begin{table}[htbp]
    \centering
    \caption{反幻觉提示词工程层级}
    \label{tbl:antihallucination}
    \begin{tabular}{lp{7cm}}
    \toprule
    层级 & 强制规则 \\
    \midrule
    引用格式 & 强制格式：[作者, 年份, 标题]。提供错误示例。禁止引用列表外论文。 \\
    反空泛表述 & "取得了较好效果"若无具体指标值则禁止使用。每个论断必须引用定量数据。 \\
    公式深度 & 每个方法必须列出核心公式，逐符号解释（含义、量纲、取值范围）。必须说明推导来源。 \\
    原理级对比 & 禁止"A优于B"而不解释公式/物理层面的原因。 \\
    理论极限 & 必须从第一性原理（物理定律、数学定理）推导性能上下界。 \\
    \bottomrule
    \end{tabular}
\end{table}

\textbf{回退机制。}
始终生成基于模板的回退报告作为安全网。若LLM失败（如3次重试后仍遇到速率限制），回退报告作为最终输出。

\textbf{语言一致性。}
系统检测研究主题的语言，在整个报告中强制使用一致的语言。中文主题生成中文报告，英文主题生成英文报告。
